\documentclass[11pt]{article}
% ======================================================================
%  Preambule commun - Sujets Bac Serie A (Mathematiques)
%  Style sobre : noir et blanc, sans couleurs, sans filigrane,
%  sans en-tetes ni pieds de page decoratifs.
% ======================================================================
\usepackage[utf8]{inputenc}
\usepackage[T1]{fontenc}
\usepackage{lmodern}
\usepackage{textcomp}
\usepackage{amsmath,amssymb}
\usepackage{enumitem}
\usepackage{array}
\usepackage[a4paper,margin=2.3cm]{geometry}
\usepackage{fancyhdr}
\usepackage{tikz}

% Symbole degre robuste + justification souple
\DeclareUnicodeCharacter{00B0}{\ensuremath{^\circ}}
\tolerance=2000
\emergencystretch=2em
\frenchspacing

% En-tete / pied de page minimalistes (numero de page seul)
\pagestyle{fancy}
\fancyhf{}
\renewcommand{\headrulewidth}{0pt}
\fancyfoot[C]{\thepage}

% Macros mathematiques
\newcommand{\R}{\mathbb{R}}
\newcommand{\N}{\mathbb{N}}
\newcommand{\Z}{\mathbb{Z}}
\newcommand{\vi}{\vec{\imath}}
\newcommand{\vj}{\vec{\jmath}}

% Titres de blocs
\newcommand{\exo}[2]{\bigskip\noindent{\large\bfseries Exercice #1}\hfill\textit{#2}\par\nopagebreak\smallskip}
\newcommand{\partie}[1]{\medskip\noindent{\bfseries Partie #1}\par\nopagebreak\smallskip}
\newcommand{\entetesujet}[1]{\begin{center}{\Large\bfseries #1}\end{center}\vspace{0.3em}\hrule\vspace{1.1em}}

% Questions numerotees : niveau 1 = chiffres gras, niveau 2 = a. b. c.
\newlist{questions}{enumerate}{1}
\setlist[questions]{label=\textbf{\arabic*.},leftmargin=1.8em,itemsep=3pt,topsep=3pt}
\newlist{sousq}{enumerate}{1}
\setlist[sousq]{label=\textbf{\alph*.},leftmargin=1.8em,itemsep=2pt,topsep=2pt}


\begin{document}

\entetesujet{Sujet bac 2015 -- Série A}

\exo{1}{8 points}
\partie{A}
\begin{questions}
  \item Définir chacune des expressions suivantes :
  \begin{itemize}
    \item Nombre rationnel.
    \item Période d'un rationnel.
  \end{itemize}
  \item Déterminer le développement décimal et la période de chacun des nombres suivants :
  \[ A = \frac{704}{37} \quad \text{et} \quad B = \frac{22}{7}. \]
  \item Déterminer l'écriture fractionnaire irréductible de chacun des nombres suivants :
  \[ C = 2{,}832 \quad \text{et} \quad E = 23{,}54 \]
\end{questions}

\partie{B}
\begin{questions}
  \item Vérifier que pour tout nombre réel $x$, on a :
  \[ (x + 1)(x^2 - 6x + 8) = x^3 - 5x^2 + 2x + 8 \]
  \item \begin{sousq}
    \item Résoudre dans $\R$ l'équation : $x^2 - 6x + 8 = 0$.
    \item Déduire de tout ce qui précède la résolution dans $\R$ de l'équation :
    \[ x^3 - 5x^2 + 2x + 8 = 0 \]
  \end{sousq}
  \item Résoudre dans $\R$ l'équation :
  \[ e^{3x} - 5\,e^{2x} + 2\,e^{x} + 8 = 0 \]
\end{questions}

\exo{2}{8 points}
On considère la fonction numérique $f$ à variable réelle $x$ définie par :
\[ f(x) = \frac{4 - x}{2x - 4} \]
$(\mathcal{C})$ désigne la courbe représentative de $f$ dans le repère $(O,\vi,\vj)$. Unité graphique : 1 cm.
\begin{questions}
  \item Déterminer l'ensemble de définition $E_f$.
  \item Calculer les limites de $f$ aux bornes de $E_f$.
  \item Calculer la dérivée $f'$ de $f$.
  \item Préciser le signe de $f'$ puis en déduire le sens de variation de $f$.
  \item Dresser le tableau de variation de $f$.
  \item Montrer que la courbe $(\mathcal{C})$ admet deux asymptotes. Préciser leurs équations.
  \item Compléter le tableau de valeurs suivant :
  \begin{center}\renewcommand{\arraystretch}{1.5}
  \begin{tabular}{|c|c|c|c|c|c|c|c|}
  \hline
  $x$ & $-3$ & $-2$ & $-1$ & $0$ & $\dfrac{3}{2}$ & $2{,}5$ & $5$ \\
  \hline
  $f(x)$ & & & & & & & \\
  \hline
  \end{tabular}
  \end{center}
  \item Construire dans le repère $(O,\vi,\vj)$ :
  \begin{itemize}
    \item les points d'abscisses : $-3$ ; $-2$ ; $-1$ ; $0$ ; $\dfrac{3}{2}$ ; $2{,}5$ ; $5$ ;
    \item les deux asymptotes ;
    \item et la courbe $(\mathcal{C})$.
  \end{itemize}
\end{questions}

\exo{3}{4 points}
Sur une même verticale, la pression atmosphérique diminue lorsque l'altitude augmente conformément au tableau suivant ($x$ : altitude en km ; $y$ : pression en cm de mercure) :
\begin{center}\renewcommand{\arraystretch}{1.4}
\begin{tabular}{|c|c|c|c|c|c|c|}
\hline
$x_i$ & 0 & 1 & 2 & 4 & 6 & 10 \\
\hline
$y_i$ & 76 & 67 & 59 & 46 & 35 & 20 \\
\hline
\end{tabular}
\end{center}
\begin{questions}
  \item Représenter graphiquement dans un repère orthonormal, par un nuage de points, la série statistique $(x_i, y_i)$.

  Unité graphique : prendre 1 cm pour 1 cm de mercure et prendre 1 cm pour 10 km en altitude.
  \item Calculer les coordonnées du point moyen $G$.
  \item Placer le point moyen $G$.
  \item On se propose de déterminer la droite d'ajustement linéaire par la méthode de Mayer.
  \begin{sousq}
    \item Découper le nuage de points en deux sous-nuages disjoints et de taille identique.
    \item Calculer les points moyens $G_1$ et $G_2$ de chaque sous-nuage.
    \item Montrer que l'équation de la droite d'ajustement est : $y = -5{,}9x + 414{,}75$.
  \end{sousq}
  \item En déduire l'altitude d'un lieu où la pression atmosphérique est de 40 cm de mercure.
\end{questions}

\end{document}
